
\documentclass[aspectratio=169]{beamer}

\mode<presentation>
\usetheme{Boadilla}
\definecolor{NTHU1}{RGB}{127,16,132}
\definecolor{NTHU2}{RGB}{228,210,231}
\definecolor{blue}{RGB}{30,90,205}
\definecolor{red}{RGB}{213,94,0}
\definecolor{green}{RGB}{0,128,0}
\definecolor{charcoalgray}{RGB}{108,104,104}
\setbeamercolor{title}{fg=NTHU1}
\setbeamercolor{frametitle}{fg=NTHU1}
\setbeamercolor{block title}{bg=NTHU1, fg=white}
\setbeamercolor{block body}{bg=NTHU2}
\setbeamercolor{structure}{fg=NTHU1}
\setbeamercolor{item projected}{fg=white}
\setbeamercolor{item}{fg=NTHU1}
\setbeamercolor{subitem}{fg=NTHU1}
\setbeamercolor{section in toc}{fg=NTHU1}
\setbeamercolor{description item}{fg=NTHU1}
\setbeamercolor{caption name}{fg=NTHU1}
\setbeamercolor{button}{bg=NTHU1, fg=white}
\setbeamercolor{caption name}{fg=NTHU1}
\usepackage{graphics}
\usepackage{tikz}
\usepackage{amsmath}
\usepackage{bbm}
\usetikzlibrary{decorations.pathreplacing}
\usepackage{geometry}
\usepackage{booktabs}
\usepackage{bookmark}
\usepackage{multirow, makecell}
\usepackage{float}
\usepackage{fancyvrb}
\usepackage{caption}
\captionsetup[figure]{singlelinecheck = false, format= hang, justification=centering}
\captionsetup[subfigure]{singlelinecheck = false, format= hang, justification=raggedright}
\usepackage{subcaption}
\usepackage{adjustbox}
\usepackage{hyperref}
\usepackage{threeparttable}
\usepackage{appendixnumberbeamer} 
\usepackage[T1]{fontenc}
\usepackage[default]{lato} 
\usepackage{smartdiagram}
\smartdiagramset{
  back arrow disabled = true,
  module minimum width=1cm,
  module x sep = 3,
  uniform arrow color=true,
}
\usepackage{xcolor}
\usepackage{colortbl}
  \definecolor{light-gray}{gray}{0.99}

\newenvironment{wideitemize}{\itemize\addtolength{\itemsep}{10pt}}{\enditemize}
\newenvironment{wideenumerate}{\enumerate\addtolength{\itemsep}{10pt}}{\endenumerate}
\newenvironment{widedescription}{\description\addtolength{\itemsep}{10pt}}{\enddescription}
\hypersetup{
colorlinks=true,
linkcolor=NTHU1,
filecolor=green, 
urlcolor=blue,
}
\beamertemplatenavigationsymbolsempty
\setbeamercolor{author in head/foot}{bg=white, fg=NTHU1}
\setbeamercolor{title in head/foot}{bg=white, fg=NTHU1}
\setbeamercolor{date in head/foot}{bg=white, fg=NTHU1}
\setbeamercolor{section in head/foot}{bg=white, fg=NTHU1}
\setbeamercolor{headline}{bg=white}
\setbeamertemplate{footline}{
    \leavevmode%
    \hbox{%
        \begin{beamercolorbox}[wd=.333333\paperwidth,ht=2.25ex,dp=1ex,center]{date in head/foot}%
            \usebeamerfont{date in head/foot}%\insertdate
        \end{beamercolorbox}%
        \begin{beamercolorbox}[wd=.333333\paperwidth,ht=2.25ex,dp=1ex,center]{title in head/foot}%
            \usebeamerfont{title in head/foot}\insertshorttitle
        \end{beamercolorbox}%
        \begin{beamercolorbox}[wd=.333333\paperwidth,ht=2.25ex,dp=1ex,right]{date in head/foot}%
            \usebeamerfont{date in head/foot}\insertframenumber{}\hspace*{2em}
        \end{beamercolorbox}}%
        \vskip0pt%
    }
    %% May need to script the introduction!!!!!!
%\setbeamercolor{page number in head/foot}{fg=NTHU1}
\setbeamertemplate{section in toc}[sections numbered]
\setbeamertemplate{subsection in toc}{\leavevmode\leftskip=3em\rlap{\hskip-1.75em\inserttocsectionnumber.\inserttocsubsectionnumber}
\inserttocsubsection\par}
\setbeamerfont{subsection in toc}{size=\footnotesize}
%\setbeamertemplate{headline}{%
  %\begin{beamercolorbox}[ht=5.5ex]{section in head/foot}
    %\vskip2pt\insertnavigation{0.33\paperwidth}\vskip2pt
  %\end{beamercolorbox}%
%}
\newenvironment{transitionframe}{\setbeamercolor{background canvas}{bg=white}\setbeamertemplate{footline}{} \begin{frame}[noframenumbering]}{\end{frame}}

\makeatletter
\let\@@magyar@captionfix\relax
\makeatother

\title[]{Conflict and Peace} % Change this regularly
\subtitle[]{Week 13: Fiscal capacity of the state}
\author[Lee]{Seung-hun Lee }
\institute[]{Taipei School of Economics\\ National Tsing Hua University}

\date[]{November 28th, 2025}

\begin{document}
\begin{transitionframe}
\titlepage
\end{transitionframe}

%%% Color slides for section headers: Use for colloquium version (The ones bewteen \iffals and \fi)

\section{Overview of the topic}
\begin{frame}
\frametitle{Development of fiscal capacity is a key feature of development} 
Expansion of government size and fiscal capacity investment in the last century
\begin{itemize}\setlength\itemsep{3pt}
\item Tax share of GDP: 10\% in 1910 $\to$ 30-50\% today
\item Tax capacity improved following implementation of broad-based taxes (income tax)
\item However, this explains the trend only in developed countries
\begin{itemize}\setlength\itemsep{3pt}
\item So economic growth $\to$ better fiscal capacity? Not regularly
\item Separate investment in capacities are key
\end{itemize}
\end{itemize}\medskip
But in poor states, no noticeable increase in fiscal capacities
\begin{itemize}\setlength\itemsep{3pt}
\item Raise significantly less revenues than developed countries
\item Lag behind in the implementation of modern tax technologies
\begin{itemize}\setlength\itemsep{3pt}
\item Property registers are manually updated, if updated at all 
\item Do not utilize third-party information that ensures accurate tax information
\item Less digitalized systems, more loopholes in general
\end{itemize}
\end{itemize}
\end{frame}

\begin{frame}
\frametitle{Rise of informational capacity and taxation over time} 
\begin{figure}
  \includegraphics[keepaspectratio, width=0.7\textwidth]{fig1.png}
\end{figure}
\end{frame}

\begin{frame}
\frametitle{Economic growth does not automatically lead to increased tax} 
\begin{figure}
  \includegraphics[keepaspectratio, width=0.7\textwidth]{fig2.png}
\end{figure}
\end{frame}

\begin{frame}
\frametitle{Less capable states rely more on distortionary and narrow based taxes} 
Why do they depend on such taxation
\begin{itemize}\setlength\itemsep{3pt}
\item Broad-base taxes require registries that cover formal sector, and audits
\item Narrow-base taxes rely on easily observable bases (trade ports, mines etc)
\item Large informal economy and underdeveloped finance sector: Opens up loopholes
\item Extractive for select groups by shifting tax burden onto them
\item Bad equilibrium: Fiscal capacity$\Downarrow\to$ Public goods and economic development $\Downarrow$
\end{itemize}\medskip
What the data tells us
\begin{itemize}\setlength\itemsep{3pt}
\item High income countries rely on broad-based taxes
\item Narrow base taxes (tariffs) more observable in low-capacity countries
\end{itemize}
\end{frame}

\begin{frame}
\frametitle{High income $\to$ Income tax, Lower income $\to$ Trade tax} 
\begin{figure}
  \includegraphics[keepaspectratio, width=0.7\textwidth]{fig3.png}
\end{figure}
\end{frame}

\begin{transitionframe}
\frametitle{Roadmap for today}
\tableofcontents
\end{transitionframe}

\section{Tax structure across developing countries}
\subsection{Gordon, Li (2009): Tax structures in developing countries: Many puzzles and a possible explanation}
\begin{transitionframe}
\begin{center}
  \begin{huge}
    \textcolor{NTHU1}{%
  Gordon, Li (2009):\\ Tax structures in developing countries: Many puzzles and a possible explanation
    }
  \end{huge}
\end{center}
\end{transitionframe}

\begin{frame}
\frametitle{Q: Do presence of informal sectors affect choice of tax policies?} 
Major difference in tax structures across developing and developed countries
\begin{itemize}\setlength\itemsep{3pt}
\item Some could be attributed to social preferences over public goods and redistribution
\item But many differences with predicted optimal taxation in generic settings
\begin{itemize}\setlength\itemsep{3pt}
\item These taxation preserves production efficiency under plausible assumptions
\item Rule out tariffs, capital income taxes, differential consumption tax, inflation tax
\end{itemize}
\item Drastic diffrences in developing countries
\begin{itemize}\setlength\itemsep{3pt}
\item Less personal income tax, more tariffs, corporate income tax, and seignorage
\end{itemize}
\end{itemize}\medskip
This paper proposes that existence of informal sector explains different tax structures
\begin{itemize}\setlength\itemsep{3pt}
\item Incorporates enforcement problems in developing countries due to informal economy
\begin{itemize}\setlength\itemsep{3pt}
\item Informal economies: Use cash over financial intermediaries
  \item Taxes on formal economy shifts activities to informal economy $\to$ Optimal tax rate $\Downarrow$
  \item Justifies inflation tax, differential tax across industries, and tariffs
\end{itemize}

\end{itemize}\medskip
\end{frame}

\begin{frame}
\frametitle{Summary of the difference in tax choices across country income} 
\begin{figure}
  \includegraphics[keepaspectratio, width=\textwidth]{gl2009_t1.png}
\end{figure}
\begin{itemize}\setlength\itemsep{3pt}
\item  High corporate income tax, tariffs, seignorage in developing countries
\item Less personal income tax
\item Statutory tax rates are similar across different countries!
\end{itemize}
\end{frame}

\begin{frame}
\frametitle{Tax policy with limited information: Differential excise taxes} 
Characterizing informal economy into a model
\begin{itemize}\setlength\itemsep{3pt}
\item Informal economy not observed by government
\begin{itemize}\setlength\itemsep{3pt}
\item Thus, taxes cannot be levied to such firms
\item So profits for such firms are revenues - wage cost - capital costs
\end{itemize}
\item Formal firms use financial systems and have higher outputs, but subject to tax
\item Firms become formal if net profits exceed those from informal economy
\begin{itemize}\setlength\itemsep{3pt}
\item High excise taxes reduce the net profits $\to$ tax revenue drops
\end{itemize}
\end{itemize}\medskip
Variation across industries
\begin{itemize}\setlength\itemsep{3pt}
\item Sectors with low increase in output from formalizing $\to$ lower tax rates than desired
\item This results in distortion in domestic production, which are offset by tariffs on imports
\item Inflation tax and red tape on lightly taxed sectors may be attractive
\end{itemize}
\end{frame}




\begin{frame}
\frametitle{How to increase benefits of using financial sector} 
Determinants of benefits to using financial sectors
\begin{itemize}\setlength\itemsep{3pt}
\item Capital intensive firms: They need lending to finance purchase of capital
\item If firms operate at wide distance, financial sector reduces the cost
\item High revenue per employee: More cash transaction exposes to internal theft
\begin{itemize}
\item Service sectors tend to be less capital intensive and easily transfer to informal economy
\end{itemize}
\item Firms that do not formalize avoid corporate and excise tax, but pay inflation tax
\item Thus, taxes are kept low for non capital-intensive sectors
\begin{itemize}\setlength\itemsep{3pt}
\item Capital-intensive ones more likely to stick to formal even without intervention
\item This is meant to shift informal economies to formal ones
\end{itemize}

\item This explains high capital income tax in developing countries (vs. income tax)
\item Inflation taxes also raise cost of keeping cash, inducing firms to formalize
\end{itemize}\medskip

\end{frame}

\begin{frame}
\frametitle{Motivating uses of other anti-competitive measures}
Tariffs  
\begin{itemize}\setlength\itemsep{3pt}
\item Higher tariff protection for capital intensive sectors
\begin{itemize}\setlength\itemsep{3pt}
\item Meant to offset distortions due to uneven domestic tax enforcement
\item Specifically, aim is to induce production to a sector more likely to formalize
\end{itemize}
\end{itemize}\medskip
Other anti-competitive measures
\begin{itemize}\setlength\itemsep{3pt}
\item Provide monopoly power to capital-intensive firms
\item Impose tighter regulation or monitoring on informal firms
\begin{itemize}\setlength\itemsep{3pt}
\item Even if revenue from such firms cannot be observed, other activities could be
\item Higher red tape measure on those firms as monitoring device
\end{itemize}
\end{itemize}\medskip
\end{frame}

\begin{frame}
\frametitle{Key takeaways}
Existance of informal firms twist the choice of optimal taxation 
\begin{itemize}\setlength\itemsep{3pt}
\item More capital taxation, tariffs, and inflation taxes
\item This is determined by the returns to using financial sectors
\begin{itemize}\setlength\itemsep{3pt}
\item If financial sectors provide benefits, less incentives to join informal economy
\item As developing countries have less developed finance sectors, high informality
\end{itemize}
\item Thus, models should incorporate possiblity of informal economies
\end{itemize}\medskip
However, leaving informal economies as is is undesirable
\begin{itemize}\setlength\itemsep{3pt}
\item Weakens capacity to provide essential public goods
\item Resource misallocation$\Uparrow$ and innovation$\Downarrow$ (cost of hiding leave little for innovation)
\item Thus, efforts to induce firms to formalize through third-party information, legal measures, and expansion of electronic filing (Benhassine et al 2018)
\end{itemize}
\end{frame}


\subsection{Okunogbe, Tourek (2024): How can lower-income countries collect more taxes? The role of technology, tax agents, and politics}
\begin{transitionframe}
\begin{center}
  \begin{huge}
    \textcolor{NTHU1}{%
  Okunogbe, Tourek (2024):\\ How can lower-income countries collect more taxes? The role of technology, tax agents, and politics
    }
  \end{huge}
\end{center}
\end{transitionframe}

\begin{frame}
\frametitle{Q: How do you improve tax system in developing countries? } 
Increasing tax revenues is a major goal in many low/middle income countries
\begin{itemize}\setlength\itemsep{3pt}
\item Generally have lower share of their GDP as tax revenues
\item Low taxation is a symptom of economic and institutional constraints
\begin{itemize}
\item Tax base not large enough, lack of constraint on redistribution/taxes
\end{itemize}
\item Does economic growth automatically expand tax capacity? Not exactly
\begin{itemize}
\item Separate investments in improving the tax system required
\end{itemize}
\end{itemize}\medskip
This paper reviews recent developments in tax administration in developing countries
\begin{itemize}\setlength\itemsep{3pt}
\item Identifying taxpayers (expanding tax registries and databases)
\item Detecting tax burden (usage of third-party information)
\item Collect taxes to state coffers (Billing systems, legal measures for noncompliance)
 \end{itemize}
\end{frame}

\begin{frame}
\frametitle{Technology in tax collection} 
Identification: Cost of linking identity info across government agencies and registries$\Downarrow$

\begin{itemize}\setlength\itemsep{3pt}
\item National ID (Ghana Card), Digital mapping of properties (Liberia)
\item ID information may need to be combined with other measures (Okunogbe 2023)
\end{itemize}\medskip
Detection: Collect third-party info (declaration of tax liabilities by third person)
\begin{itemize}\setlength\itemsep{3pt}
\item `Digital trails' from value-added taxes and employment taxes
\item Requires countries to cross-check information using multiple sources
\item Need to close international loopholes
\end{itemize}\medskip
Collection: Increases in electronic filing and mobile payment options
\begin{itemize}\setlength\itemsep{3pt}
\item Reduces cost of monitoring and Compliance
\item May disadvantage those less familiar with electronics \\ (rural, less-educated, less-established firms)
\end{itemize}\medskip

\end{frame}

\begin{frame}
\frametitle{Cross-country evidence on technology and collection} 
\begin{figure}
  \includegraphics[keepaspectratio, width=0.6\textwidth]{ot2024_f1.png}
\end{figure}
\end{frame}

\begin{frame}
\frametitle{Role of tax collectors} 
What determines the effectiveness of officials
\begin{itemize}\setlength\itemsep{3pt}
\item Deployment: low citizen-to-staff ratio and appropriate allocation
\begin{itemize}\setlength\itemsep{3pt}
\item Low burden per staff
\item Specialization: Matching staff depending on tax types and taxpayer traits 
\end{itemize}
\item Incentives: Tie staff revenues to tax revenue generations
 \begin{itemize}\setlength\itemsep{3pt}
\item Risk: Could increase scope for corruption
\item Monitoring staff with audits crucial 
\end{itemize}
\end{itemize}\medskip
Interaction with technology
\begin{itemize}\setlength\itemsep{3pt}
\item Substitute? Obviate need for in-person Interaction
\item Complement? Help expedite tedious processes and collecting contextual knowledge
\item Balancing the two forces is key (and unresolved question)
\item Resistance to adopting the technological change
\end{itemize}\medskip
\end{frame}


\begin{frame}
\frametitle{Political incentives for taxation} 
How do political incentives determine tax collection
\begin{itemize}\setlength\itemsep{3pt}
\item Tax$\Uparrow$ $\to$ fund public services, social programs, infrastructure $\to$ support re-election
\item Tax$\Uparrow$ $\to$ Citizen demand, expectation for accountability $\Uparrow$ $\to$ not help that much
\end{itemize} \medskip
What are the factors shaping success of tax collection?
\begin{itemize}
\item Alternative revenue: Dampen motivation for Taxation
\item Political competition: Competition for better governance (through public services)
\item Availability of tech: Reduce cost of collection and resistance from taxpayers
\end{itemize}

\end{frame}

\begin{frame}
\frametitle{Expanding our knowledge} 
Findings from other relevant studies
\begin{itemize}\setlength\itemsep{3pt}
\item Tax revenue and compliance increases with increased tech and enforcement
\item Many results are from small-scale local experiments: Generalizable?
\end{itemize}\medskip
What remains?
\begin{itemize}\setlength\itemsep{3pt}
\item How to make technology and personnel complementary?
\item How to tax new commodities (digital activities, digital currencies)?
\item How to increase political support for taxation?
\end{itemize}
\end{frame}

\subsection{Jensen (2022): Employment Structure and the Rise of the Modern Tax System}
\begin{transitionframe}
\begin{center}
  \begin{huge}
    \textcolor{NTHU1}{%
  Jensen (2022):\\ Employment Structure and the Rise of the Modern Tax System
    }
  \end{huge}
\end{center}
\end{transitionframe}

\begin{frame}
\frametitle{Q: How do modern tax systems rise?}
Modern tax systems are characterized by increased direct taxes (income tax)
\begin{itemize}\setlength\itemsep{3pt}
\item Requires high collection capacities 
\item Changes in the structure of employment may explain increase in collection
\begin{itemize}\setlength\itemsep{3pt}
\item Enforcement in employer-employee structure $>$ Self-employed
\item Existence of third-party information (employer, in case of income tax)
\end{itemize}
\end{itemize}\medskip
This paper uncovers new facts on tax capacity with microdata across countries and time
\begin{itemize}\setlength\itemsep{3pt}
\item Employee share increases over income distribution \& further down the income distribution as country develops 
\item Exemption threshold gradually decreases in income distribution over development
\item Threshold is located such that employee share on tax base remains constant and high
\end{itemize}
\end{frame}

\begin{frame}
\frametitle{Data, definition, and calculation strategy} 
Data source and definitions
\begin{itemize}\setlength\itemsep{3pt}
\item Household surveys (100 countries)
\begin{itemize}\setlength\itemsep{3pt}
\item Employment type of economically active population (EAP)
\item Information on taxable earnings (employment / self-employment / capital / other income)
\item For US (1870-2010) and Mexico (1960-2010), compile long run time series 
\end{itemize}
\end{itemize}\medskip
Definition of key terms and strategies for calculation
\begin{itemize}\setlength\itemsep{3pt}
\item Employee: if work type generates third-party information usable in tax enforcement \\ (Otherwise, self-employment)
\begin{itemize}\setlength\itemsep{3pt}
\item Self-employed: casual laborers, own account workers, workers w/o contracts
\item Calculate share of employees in EAP across country's income distribution (in deciles)
\end{itemize}
\item Threshold: Minimum level of gross income above which individuals have to pay taxes
\begin{itemize}\setlength\itemsep{3pt}
\item From country tax reports of International Bureau of Fiscal Documentation
\end{itemize}
\item Income tax base: \% EAP whose gross income is above income tax exemption threshold

\end{itemize}
\end{frame}

\begin{frame}
\frametitle{Fact1: Employee share $\Uparrow$ with income and development levels} 
\begin{columns}
\begin{column}{0.6\textwidth}
  \begin{figure}
    \includegraphics[keepaspectratio, width=\textwidth]{j2022_f1.png}
    \caption*{Red: Employee, Blue: Self-employed}
  \end{figure}
  
\end{column}  
\begin{column}{0.35\textwidth}
\begin{itemize}\setlength\itemsep{3pt}
\item Within country, employee share increase over income
\item Across countries, employee share profile moves leftward 
\item Over transition, increase in employee share concentrated in deciles gradually further down 
\end{itemize}  
\end{column}
\end{columns}
\end{frame}

\begin{frame}
\frametitle{Fact2: Exemption threshold moves downward as country develops} 
\begin{columns}
\begin{column}{0.6\textwidth}
  \begin{figure}
    \includegraphics[keepaspectratio, width=\textwidth]{j2022_f2.png}
    \caption*{Red: Employee, Blue: Self-employed, Vertical: Threshold}
  \end{figure}
  
\end{column}  
\begin{column}{0.35\textwidth}
\begin{itemize}\setlength\itemsep{3pt}
\item At low development, threshold at top deciles 
\item With development, growth in employee share comoves with falling thresholds
\item Tax base increases
\item Documents frequent updating of thresholds (base increase not driven by nominal inflation)
\end{itemize}  
\end{column}
\end{columns}
\end{frame}

\begin{frame}
\frametitle{Fact3: Threshold is located where employee share remains constant} 
\begin{columns}
\begin{column}{0.6\textwidth}
  \begin{figure}
    \includegraphics[keepaspectratio, width=\textwidth]{j2022_f3.png}
    \caption*{}
  \end{figure}
  
\end{column}  
\begin{column}{0.35\textwidth}
\begin{itemize}\setlength\itemsep{3pt}
\item Composition of tax base is constant and high in employee share
\item Employee share: 85-95\%
\item Suggest that employee composition matters in deciding threshold
\item \% Self-employed $\Uparrow\to$ Enforcement costs $\Uparrow$ (high employee share desirable)
\end{itemize}  
\end{column}
\end{columns}
\end{frame}

\begin{frame}
\frametitle{Takeaways from the stylized facts}
Cost channels at play in expanding tax base
\begin{itemize}\setlength\itemsep{3pt}
\item Concerns with taxing and income segment with too high-self employed population
\begin{itemize}
\item Employees have information trails, self-employed do not
\end{itemize}
\item Compliance cost imposed on newly taxed self-employed may be high
\begin{itemize}
\item Employees have third-party withholding
\end{itemize}
\item Worsen between-group horizontal equity among employee vs self-employed
\begin{itemize}
\item Self-employed have evasion opportunities $\to$ effective taxes paid < legal liabilities 
\end{itemize}
\end{itemize}\medskip
Rise of modern taxation: Decrease in threshold (broad-based tax) happens when
\begin{itemize}\setlength\itemsep{3pt}
\item Increase in employee share at lower end of income distribution (Fact 1)
\item This lowers costs of reform and expand tax base (Fact 2)
\item While retaining high employee share, making tax administration manageable (Fact 3)
\end{itemize}
\end{frame}

\begin{frame}
\frametitle{Implications for developing countries} 
Redistribution by direct taxes might be ill-suited for developing countries
\begin{itemize}\setlength\itemsep{3pt}
\item Developed countries extensively use third-party information
\begin{itemize}
\item Evidenced by increase in employee share down the income distribution over time
\end{itemize}
\item Developing countries cannot enforce taxes flexibly due to cost constraints 
\begin{itemize}
\item Higher share of self-employment (difficult to gather info)
\end{itemize}
\item For such settings, indirect taxation could be the best option for redistribution
\end{itemize}\medskip
Remaining issues 
\begin{itemize}
\item How would implementation of technology affect tax base long-run?
\begin{itemize}\setlength\itemsep{3pt}
\item More use of third-party information $\to$ Rising `employee' share?
\item Expanding tax base would then be less costly $\to$ tax collection, public goods $\Uparrow$
\item But what about adoption (existence of political will)?
\end{itemize}
\end{itemize}
\end{frame}

\subsection{Olsson, Baaz, and Martinsson (2020): Fiscal capacity in post-conflict states: Evidence from trade on Congo River [student presentation]}
\begin{transitionframe}
\begin{center}
  \begin{huge}
    \textcolor{NTHU1}{%
Olsson, Baaz, and Martinsson (2020):\\ Fiscal capacity in post-conflict states: Evidence from trade on Congo River [student presentation]
    }
  \end{huge}
\end{center}
\end{transitionframe}

\begin{frame}
\frametitle{Course logistics} 
\begin{itemize}\setlength\itemsep{3pt}
\item 
\begin{itemize}
\item 
\end{itemize}
\end{itemize}
\end{frame}


\section{Barriers to taxation in developing countries}
\subsection{Balán, Bergeron, Tourek, Weigel (2022): Local Elites as State Capacity: How City Chiefs Use Local Information to Increase Tax Compliance in the DRC}
\begin{transitionframe}
\begin{center}
  \begin{huge}
    \textcolor{NTHU1}{%
Balán, Bergeron, Tourek, Weigel (2022):\\ Local Elites as State Capacity: How City Chiefs Use Local Information to Increase Tax Compliance in the DRC
    }
  \end{huge}
\end{center}
\end{transitionframe}

\begin{frame}
\frametitle{Q: How can local information be used effectively to increase taxes?} 
How can fiscal capacities develop in presence of local \& traditional elites
\begin{itemize}\setlength\itemsep{3pt}
\item On the one hand, can capture politics and civic society
\item On the other, could be room to collaborate to improve governance
\item Potential trade-offs
\begin{itemize}
\item Enforcement backed by the state vs local knowledge and coordination?
\end{itemize}
\end{itemize}\medskip
This study investigates if local elites can contribute to statebuilding
\begin{itemize}\setlength\itemsep{3pt}
\item Comparison between state agents vs local elites
\begin{itemize}\setlength\itemsep{3pt}
\item Local elites know more contextual knowledge but could lead to leakage in revenues
\item State agents have better enforcement capacity
\end{itemize}
\item Based on field experiment randomizing type of tax collectors across neighborhoods
\item Estimate changes in revenues and bribes, with tests for various mechanism
\end{itemize}
\end{frame}

\begin{frame}
\frametitle{Taxation in Democratic Republic of Congo (DRC)} 
Setting of DRC
\begin{itemize}\setlength\itemsep{3pt}
\item Low-capacity state: tax/GDP ranks 188 out of 200 and  \$0.30 per capita tax revenue
\item Relies on property tax: First campaign to expand property tax in 2016
\item Public goods and services are scarce and low-quality
\begin{itemize}
\item Limited running water, electricity, unrepaired roads etc
\end{itemize}
\end{itemize}\medskip
The current study focuses on the second campaign (May-Dec 2018)
\begin{itemize}\setlength\itemsep{3pt}
\item Collectors trained by provincial ministry, separate for state and chiefs
\item Involved property registration updates, providing liability information 
\item Followed by tax visits (revisit until taxes are paid)
\item Collectors compensated based on houses registered and revenue collection
\item Property tax rates are tiered flat, fixed fees (low: \$2 tax, 89\% vs high: \$9 tax, 11\%)
\begin{itemize}
\item Equivalent to 0.32\% of property value ($\simeq$ lower end of property tax rates in US)
\end{itemize}
\end{itemize}\medskip
\end{frame}

\begin{frame}
\frametitle{Data and empirical strategy} 
RCT on 45000+ properties in 356 neighborhoods in Kasai-Central
\begin{itemize}\setlength\itemsep{3pt}
\item Randomized at neighborhood, treatment at past campaigns, experience of chiefs
\item Central (state agents) vs Local (chiefs) vs State agents w/ local info vs state and local
\item Pure control group: Declare their property values voluntarily
\item Other aspects of collection equal
\begin{itemize}
\item Property registration, assessment, tax liabilities, training, collection compensation
\end{itemize}
\item Complemented with tax administrative data and household surveys
\end{itemize}\medskip
Estimation leverages random assignment of treatment (indiv $i$, neighbor $j$, strata $k$, time $t$)
\[
y_{ijkt}=\beta_0 + \beta\text{Local}_{jkt}+\gamma X_{ijkt} + \alpha_k + \theta_t + e_{ijkt}
\]\vspace{-20pt}
\begin{itemize}\setlength\itemsep{3pt}
\item For mechanisms, other treatment arms included as separate variable
\end{itemize}
\end{frame}

\begin{frame}
\frametitle{Increase in compliance and revenues} 
\begin{figure}
  \includegraphics[keepaspectratio, width=0.8\textwidth]{bbtw2022_t1.png}
\end{figure}
\end{frame}

\begin{frame}
\frametitle{Accurate assessments \& trust in government $\uparrow$, but bribes $\uparrow$} 
\begin{figure}
  \includegraphics[keepaspectratio, width=0.7\textwidth]{bbtw2022_t2.png}
\end{figure}
\end{frame}

\begin{frame}
\frametitle{Evidence of local information at play} 
\begin{figure}
  \includegraphics[keepaspectratio, width=1\textwidth]{bbtw2022_t3.png}
\end{figure}
\end{frame}

\begin{frame}
\frametitle{Discussion and takeaways}
What explains the chiefs' advantage? Better targeting with information advantage
\begin{itemize}\setlength\itemsep{3pt}
\item Lower effort cost? Chiefs did not make more visits
\item Target their visits? When state agents received info from chiefs, compliance $\Uparrow$
\item Better persuasion? Neutralize info advantage (visit in linear route)$\Rightarrow$ Central = Local
\end{itemize}\medskip
Takeaways and remaining questions
\begin{itemize}\setlength\itemsep{3pt}
\item Should countries with low capacities rely on local elites?
\begin{itemize}\setlength\itemsep{3pt}
\item Raise revenues, but also bribes
\item To shift to state, value of \$1 lost in bribes $>$$>$ \$1 gained in revenues
\end{itemize}
\item For technology to increase taxes, need to surpass some state capacity threshold?
\begin{itemize}\setlength\itemsep{3pt}
\item Speak to the complementary/substitutional relationship between tech and person
\end{itemize}
\end{itemize}\medskip
\end{frame}


\subsection{Henn, Mugaruka, Ortiz, Sánchez de la Sierra, Wu (2025): Monopoly of Taxation without a Monopoly of Violence: The Weak State's Trade-off from Taxation}
\begin{transitionframe}
\begin{center}
  \begin{huge}
    \textcolor{NTHU1}{%
  Henn, Mugaruka, Ortiz, Sánchez de la Sierra, Wu (2025):\\ Monopoly of Taxation without a Monopoly of Violence: The Weak State's Trade-off from Taxation
    }
  \end{huge}
\end{center}
\end{transitionframe}

\begin{frame}
\frametitle{Q: Tradeoffs between monopoly of taxation \& monopoly of violence?} 
Weak states often fail to exert de facto control over their jurisdictions
\begin{itemize}\setlength\itemsep{3pt}
\item Non-state actors (rebels) take over, who sometimes provide stability over 
\item If these areas are economically valuable, state loses lot of fiscal revenue
\item So is it beneficial for state to take formal authority over those lands?
\begin{itemize}\setlength\itemsep{3pt}
\item Regain lost fiscal revenues
\item But could trigger inadvertent response from rebels (e.g. plundering)
\end{itemize}
\end{itemize}\medskip
This paper investigates the effect of state-building attempts to increase tax revenues
\begin{itemize}\setlength\itemsep{3pt}
\item Military powers to reclaim land and replace rebels
\item Negotiate with rebels (often with concession to the rebels)
\item Leverage subnational evidence in Democratic Republic of Congo (DRC) 
\end{itemize}
\end{frame}

\begin{frame}
\frametitle{State-building efforts to drive rebels out of DRC} 
State absence in Eastern DRC (and postcolonial Africa in general)
\begin{itemize}\setlength\itemsep{3pt}
\item Many parts of DRC (East) were under armed group control since 1998
\item Most notable of these organizations is FDLR (Forces de Libération du Rwanda)
\end{itemize}\medskip
Military operation and peace treaty with rebels
\begin{itemize}\setlength\itemsep{3pt}
\item Largest peace treaty: Sun City peace treaty in 2004
\begin{itemize}\setlength\itemsep{3pt}
\item Various armed groups integrated into Congolese army
\item Allowed government to regain half the country without triggering violence
\item Other peace agreements had smaller goals and lower credibility
\end{itemize}
\item Largest military operation: Kimia II (March 2009)
\begin{itemize}\setlength\itemsep{3pt}
\item Joint operation with UN, national army, and Rwanda
\item Targeted FDLR presiding and taxing over Basile Chiefdom region since 2005
\item Although defeated, FDLR launched theft operations from the nearby hilly forest region
\end{itemize}
\end{itemize}
\end{frame}

\begin{frame}
\frametitle{Data and empirical strategy} 
Constructing the data
\begin{itemize}\setlength\itemsep{3pt}
\item Armed group presence: village-level data through interviews with former combatants
\begin{itemize}
\item Perpetrators' affiliation, motivation (pilage/retaliate/etc), actions taken
\end{itemize}
\end{itemize}\medskip
Effects of Kimia II and Sun City Peace Treaty
\begin{itemize}\setlength\itemsep{3pt}
\item Kimia II: FDLR vs non-FDLR villages pre (`05-`09) and post (`10-`12) in Basile Chiefdom
\item Peace treaty: Integrated villages vs others in the region before and after `02
\[
Y_{i,t} = \alpha_i + \beta_t + \sum_{k=-4}^{k=3}\beta_k \text{Treat}_i \times 1[t=t_0+k] + e_{i,t}
\]\vspace{-20pt}
\begin{itemize}
  \item $t_0=2009$ for Kimia II effects, $2002$ for peace treaty effects
\end{itemize}
\end{itemize}
\end{frame}

\begin{frame}
\frametitle{Kimia II: FDLR increased violence afterwards} 
\begin{figure}
\includegraphics[keepaspectratio,width=0.9\textwidth]{hmosw2025_f1.png}  
\end{figure}
\end{frame}

\begin{frame}
\frametitle{Kimia II: Driven by increase in pillages by FDLR} 
\begin{figure}
\includegraphics[keepaspectratio,width=0.9\textwidth]{hmosw2025_t1.png}  
\end{figure}
\end{frame}

\begin{frame}
\frametitle{Peace treaty: Rise of new armed groups since treaty} 
\begin{figure}
\includegraphics[keepaspectratio,width=0.7\textwidth]{hmosw2025_f2.png}  
\end{figure}
\end{frame}

\begin{frame}
\frametitle{Takeaways and discussions} 
What (doesn't) explain the effects
\begin{itemize}\setlength\itemsep{3pt}
\item Actual violence increased (as opposed to \textit{reported} violence)
\item Violence by other actors do not increase following Kimia II
\begin{itemize}
\item Thus when incentive to protect villages (for tax revenues) vanish, violence increase
\end{itemize}
\item Treaty raised revenues, but limited commitment ability led to questioning of legitimacy
\end{itemize}\medskip
Takeaways and questions
\begin{itemize}\setlength\itemsep{3pt}
\item Kimia II: Monopolize violence (rebels out), but lost taxation due to ensuing violence
\item Peace treaty: Expand tax revenues but no longer monopolize violence
\item If faced with these trade-offs what should the states do?
\begin{itemize}\setlength\itemsep{3pt}
  \item Asserting force desirable if state has more capacity to prevent crime 
  \item Bargaining potentially desirable if states can commit to not negotiate with future rebels and have resource to integrate rebels
\end{itemize}
\end{itemize}
\end{frame}



\subsection{Ch, Shapiro, Steel, Vargas (2018): Endogenous Taxation in Ongoing Internal Conflict: The Case of Colombia [student presentation]}
\begin{transitionframe}
\begin{center}
  \begin{huge}
    \textcolor{NTHU1}{%
  Ch, Shapiro, Steel, Vargas (2018): \\ Endogenous Taxation in Ongoing Internal Conflict: The Case of Colombia [student presentation]
    }
  \end{huge}
\end{center}
\end{transitionframe}

\begin{frame}
\frametitle{Takeaways and remaining questions} 
Trajectory of fiscal capacity development
\begin{itemize}\setlength\itemsep{3pt}
\item More broad-based taxation (income tax) aided by investments in capacity
\item Many barriers to investment in developing countries
\begin{itemize}\setlength\itemsep{3pt}
\item Informal economies that distort taxation decisions
\item High enforcement cost due to lack of technology, third-party info, and personnel
\item Establishing legitimacy difficult in face of internal conflicts
\end{itemize}
\item Utilization of local knowledge possible: But need to be wary of local capture/bribes
\end{itemize}\medskip
Remaining questions
\begin{itemize}\setlength\itemsep{3pt}
\item Other input in state capacity: Personnel to execute tasks
\begin{itemize}\setlength\itemsep{3pt}
\item Responsible for collecting taxes, delivering public services
\item How do you recruit talented, motivated bureaucrats? How do you retain them?
\item How do you incentivize bureaucrats and monitor them to prevent deviant actions?
\end{itemize}
\end{itemize}
\end{frame}


%%%%%%%%%%%
\end{document}


